\documentclass[11pt]{article}
\usepackage[margin=1in]{geometry}
\usepackage[T1]{fontenc}
\usepackage{lmodern,microtype,hyperref,enumitem,amsmath,amssymb,xcolor,booktabs}
\usepackage[numbers,sort&compress]{natbib}
\hypersetup{colorlinks=true,urlcolor=blue,citecolor=blue,linkcolor=blue}
\title{Q18: Comprehension Check}
\author{DLS 2026 --- Kukla}
\date{October 2026}
\begin{document}
\maketitle
\textbf{Comprehension check.} A researcher shows that a language model solves a reasoning task with the
same accuracy and error pattern as humans, and concludes ``the model thinks like a brain.'' Name the
specific inference error, and state what \emph{additional} kind of evidence would be needed to support a
claim of mechanistic similarity.

\vspace{0.8cm}\hrule\vspace{4pt}
{\small\textbf{Answer.} The error is inferring \emph{mechanistic} similarity from \emph{behavioral}
similarity---matching outputs (and even error patterns) does not show the internal process is the same.
Supporting mechanistic similarity would require evidence about the computation itself: e.g.\ that the
model's internal representations predict neural activity during the same task, that identified circuits
correspond to known neural processes, or that causal interventions (ablation/patching) produce
effects analogous to lesion/stimulation studies in brains. Representational and causal-mechanistic
evidence, not behavioral matching, is what the claim needs.}
\end{document}
